\chapter[Lezione 17 --- 30 ottobre 2026]{Lezione 17: successioni e limiti notevoli}
\label{chap:lezione-17}

\section{Successioni}
\label{sec:lez17-successioni}

\begin{definition}[Successione reale]
\label{def:lez17-successione}
Sia \(A\subseteq\N\) non vuoto. Una \emph{successione reale}
è un'applicazione
\[
  a\colon A\longrightarrow\R,
  \qquad n\longmapsto a(n).
\]
Il valore \(a(n)\) si indica anche con \(a_n\), e la successione
si denota con \((a_n)_{n\in A}\) oppure con
\(\{a_n\}_{n\in A}\).
\end{definition}

\begin{example}[Successione dei reciproci]
\label{ex:lez17-successione-reciproci}
La funzione
\[
  a\colon\N\setminus\{0\}\longrightarrow\R,
  \qquad n\longmapsto\frac1n,
\]
definisce la successione \(a_n=1/n\), per \(n\ge1\).
\end{example}

L'insieme \(\N\), considerato come sottoinsieme di \(\R\),
non ha punti di accumulazione reali: per ogni \(x_0\in\R\)
esiste \(\delta>0\) tale che
\[
  ((x_0-\delta,x_0+\delta)\cap\N)\setminus\{x_0\}
  =\varnothing.
\]
Inoltre, \(\N\) è inferiormente limitato e non è
superiormente limitato. Pertanto, nella retta reale estesa,
l'unico suo punto di accumulazione è \(+\infty\).
Lo stesso vale per \(\N\setminus\{0\}\).

Per una successione definita su un sottoinsieme non limitato
di \(\N\), la nozione di limite è dunque quella della
\cref{def:lez08-limite} per \(n\to+\infty\).

\begin{example}[Limite della successione \(1/n\)]
\label{ex:lez17-limite-reciproci}
Si verifica che
\[
  \lim_{n\to+\infty}\frac1n=0.
\]
Sia \(\varepsilon>0\). Per ogni
\(n\in\N\setminus\{0\}\), si ha
\[
  \left|\frac1n-0\right|<\varepsilon
  \quad\Longleftrightarrow\quad
  n>\frac1\varepsilon.
\]
Scegliendo \(K_\varepsilon:=1/\varepsilon\), risulta
\[
  \left|\frac1n\right|<\varepsilon
  \qquad\forall n\in(K_\varepsilon,+\infty)
  \cap(\N\setminus\{0\}).
\]
Per la \cref{def:lez08-limite}, segue la tesi.
\end{example}

\clearpage

\section{Teorema ponte}
\label{sec:lez17-teorema-ponte}

Il limite di una funzione può essere caratterizzato
mediante i valori assunti lungo successioni di punti
del suo dominio che tendono al punto considerato.

\begin{theorem}[Teorema ponte o criterio successionale per i limiti]
\label{thm:lez17-teorema-ponte}
Sia \(f\colon A\longrightarrow\R\), con
\(A\subseteq\R\) non vuoto, e sia
\(x_0\in\overline{\R}\) un punto di accumulazione per \(A\).
Per ogni \(\ell\in\overline{\R}\), vale
\[
  \lim_{x\to x_0}f(x)=\ell
  \quad\Longleftrightarrow\quad
  \lim_{n\to+\infty}f(a_n)=\ell,
\]
dove la condizione a destra deve valere per ogni successione
\((a_n)_{n\ge1}\) a valori in \(A\setminus\{x_0\}\)
tale che
\[
  \lim_{n\to+\infty}a_n=x_0.
\]
\end{theorem}

\begin{proof}[Dimostrazione dell'implicazione \((\Rightarrow)\)]
Si supponga che
\[
  \lim_{x\to x_0}f(x)=\ell.
\]
Sia \((a_n)_{n\ge1}\) una successione a valori in
\(A\setminus\{x_0\}\) tale che \(a_n\to x_0\).
Applicando il \cref{thm:lez16-limite-sostituzione}
con \(\varphi(n):=a_n\) e \(t_0=+\infty\),
la prima condizione del teorema è soddisfatta,
poiché \(a_n\ne x_0\) per ogni \(n\ge1\).
Pertanto
\[
  \lim_{n\to+\infty}f(a_n)=\ell.
\]
\end{proof}

La dimostrazione dell'implicazione \((\Leftarrow)\)
non viene sviluppata in questa sede.

Il \cref{thm:lez17-teorema-ponte} permette, in particolare,
di escludere l'esistenza di un limite trovando due successioni
che tendono allo stesso punto, ma lungo le quali la funzione
ammette limiti distinti.

\begin{example}[Non esistenza del limite del seno all'infinito]
\label{ex:lez17-seno-infinito}
Si considerino le successioni
\[
  a_n:=\frac\pi2+2\pi n,
  \qquad b_n:=\pi n,
  \qquad n\ge1.
\]
Entrambe tendono a \(+\infty\), ma
\[
  \sin(a_n)=1,
  \qquad \sin(b_n)=0
  \qquad\forall n\ge1.
\]
Quindi
\[
  \lim_{n\to+\infty}\sin(a_n)=1,
  \qquad
  \lim_{n\to+\infty}\sin(b_n)=0.
\]
Se esistesse \(\lim_{x\to+\infty}\sin x=\ell\),
per il \cref{thm:lez17-teorema-ponte} entrambi i
limiti precedenti dovrebbero essere uguali a \(\ell\),
contrariamente a quanto ottenuto. Pertanto
\[
  \nexists\lim_{x\to+\infty}\sin x.
\]
\end{example}

\clearpage

\section{Limiti notevoli}
\label{sec:lez17-limiti-notevoli}

\begin{theorem}[Numero di Nepero]
\label{thm:lez17-numero-nepero}
Esiste il limite
\[
  \lim_{n\to+\infty}\left(1+\frac1n\right)^n=e,
  \qquad 2<e<3.
\]
La costante \(e\) è detta \emph{numero di Nepero}.
\end{theorem}

La dimostrazione del \cref{thm:lez17-numero-nepero}
non viene sviluppata in questa sede.

\begin{theorem}[Limiti notevoli]
\label{thm:lez17-limiti-notevoli}
Valgono i seguenti limiti:
\begin{enumerate}
  \item
  \[
    \lim_{x\to+\infty}\left(1+\frac1x\right)^x=e;
  \]
  \item
  \[
    \lim_{x\to-\infty}\left(1+\frac1x\right)^x=e;
  \]
  \item
  \[
    \lim_{x\to0}(1+x)^{1/x}=e;
  \]
  \item per \(a>0\), con \(a\ne1\),
  \[
    \lim_{x\to0}\frac{\log_a(1+x)}{x}=\log_a e;
  \]
  \item per \(a>0\), con \(a\ne1\),
  \[
    \lim_{x\to0}\frac{a^x-1}{x}=\log_e a;
  \]
  \item per \(\alpha\in\R\),
  \[
    \lim_{x\to0}\frac{(1+x)^\alpha-1}{x}=\alpha;
  \]
  \item
  \[
    \lim_{x\to0}\frac{\sin x}{x}=1;
  \]
  \item
  \[
    \lim_{x\to0}\frac{1-\cos x}{x^2}=\frac12;
  \]
  \item per \(x\in\R\),
  \[
    \lim_{n\to+\infty}\left(1+\frac{x}{n}\right)^n=e^x;
  \]
  \item per \(\alpha\in\R\),
  \[
    \lim_{x\to+\infty}\left(1+\frac\alpha x\right)^x=e^\alpha;
  \]
  \item per \(\alpha\in\R\),
  \[
    \lim_{x\to-\infty}\left(1+\frac\alpha x\right)^x=e^\alpha.
  \]
\end{enumerate}
\end{theorem}

\begin{proof}[Dimostrazione del quarto punto]
Siano \(a>0\), con \(a\ne1\), e \(x>-1\), con
\(x\ne0\). Per le proprietà dei logaritmi si ha
\[
  \frac{\log_a(1+x)}{x}
  =\log_a\left((1+x)^{1/x}\right).
\]
Per il terzo punto del \cref{thm:lez17-limiti-notevoli},
\[
  \lim_{x\to0}(1+x)^{1/x}=e.
\]
Si può applicare il \cref{thm:lez16-limite-sostituzione},
considerando il logaritmo in \(e\). Si verifica infatti che
\[
  \lim_{u\to e}\log_a u=\log_a e.
\]
Per giustificare quest'ultima uguaglianza, fissato
\(\varepsilon>0\), si considerino i due numeri
\(a^{\log_a e-\varepsilon}\) e
\(a^{\log_a e+\varepsilon}\), che si trovano da parti
opposte rispetto a \(e\). Esiste quindi un intorno
di \(e\) contenuto nell'intervallo avente tali estremi.
Per la monotonia del logaritmo, in questo intorno
\[
  \log_a e-\varepsilon<\log_a u
  <\log_a e+\varepsilon.
\]
Inoltre, l'uguaglianza \(\log_a(e)=\log_a e\)
assicura la seconda condizione del teorema di sostituzione.
Pertanto
\[
  \lim_{x\to0}\frac{\log_a(1+x)}{x}
  =\lim_{u\to e}\log_a u=\log_a e.
\]
\end{proof}

\begin{proof}[Dimostrazione del quinto punto]
Sia \(a>0\), con \(a\ne1\), e si ponga
\[
  t:=a^x-1.
\]
Per \(x\to0\) si ha \(t\to0\). Tale limite segue
dalla monotonia dell'esponenziale e dalla
\cref{prop:lez17-radici-en}: infatti, per \(a>1\),
se \(|x|<1/n\), allora
\[
  a^{-1/n}\le a^x\le a^{1/n},
\]
e i due estremi tendono a \(1\) per \(n\to+\infty\).
Il caso \(0<a<1\) si ricava ponendo \(c:=1/a>1\)
e usando \(a^x=1/c^x\).

Poiché \(a^x\ne1\) per \(x\ne0\), si ha anche
\(t\ne0\), e dalla relazione \(a^x=1+t\) segue
\[
  x=\log_a(1+t).
\]
Poiché \(\log_a e\ne0\), dal quarto punto già dimostrato
e dal \cref{thm:lez14-limite-rapporto} segue
\[
  \lim_{t\to0}\frac{t}{\log_a(1+t)}
  =\frac1{\log_a e}=\log_e a,
\]
dove l'ultima uguaglianza segue dalla formula di
cambiamento di base dei logaritmi.
Applicando il \cref{thm:lez16-limite-sostituzione},
si conclude
\[
  \lim_{x\to0}\frac{a^x-1}{x}
  =\lim_{t\to0}\frac{t}{\log_a(1+t)}
  =\log_e a.
\]
\end{proof}

Le dimostrazioni degli altri limiti del
\cref{thm:lez17-limiti-notevoli} saranno sviluppate
successivamente nel corso.

\clearpage

\section{Limiti di successioni notevoli}
\label{sec:lez17-limiti-successioni-notevoli}

\begin{proposition}[Limite del fattoriale]
\label{prop:lez17-limite-fattoriale}
Si ha
\[
  \lim_{n\to+\infty}n!=+\infty.
\]
\end{proposition}

\begin{proof}
Per ogni \(n\ge1\), si ha \(n!\ge n\).
Poiché \(n\to+\infty\), il
\cref{thm:lez12-confronto-secondo} permette di concludere
che
\[
  \lim_{n\to+\infty}n!=+\infty.
\]
\end{proof}

\begin{proposition}[Limite della radice \(n\)-esima]
\label{prop:lez17-radici-en}
Per ogni \(a>0\), si ha
\[
  \lim_{n\to+\infty}a^{1/n}=1.
\]
\end{proposition}

\begin{proof}
Si distinguono tre casi.

\emph{Primo caso: \(a=1\).}
Poiché \(a^{1/n}=1\) per ogni \(n\ge1\),
la tesi è immediata.

\emph{Secondo caso: \(a>1\).}
Si ponga
\[
  b_n:=a^{1/n}-1>0.
\]
Per la \cref{thm:disuguaglianza-bernoulli},
per ogni \(n\ge1\), si ha
\[
  a=(1+b_n)^n\ge1+nb_n,
\]
da cui
\[
  0<b_n\le\frac{a-1}{n}.
\]
Poiché \(a>1\) è fissato e
\[
  \lim_{n\to+\infty}\frac{a-1}{n}=0,
\]
per il \cref{thm:lez12-tre-funzioni} si ottiene
\[
  \lim_{n\to+\infty}b_n=0.
\]
Essendo \(a^{1/n}=1+b_n\), per il
\cref{thm:lez14-limite-somma} si conclude
\[
  \lim_{n\to+\infty}a^{1/n}=1.
\]

\emph{Terzo caso: \(0<a<1\).}
Si ponga \(c:=1/a>1\). Per il secondo caso,
\[
  \lim_{n\to+\infty}c^{1/n}=1.
\]
Poiché \(a^{1/n}=1/c^{1/n}\), per il
\cref{thm:lez14-limite-rapporto} segue
\[
  \lim_{n\to+\infty}a^{1/n}
  =\lim_{n\to+\infty}\frac1{c^{1/n}}=1.
\]
\end{proof}
